\section{Proofs for retracts and the constructed Plotkin powerdomain}
\label{annex:p10-domain-constructions-retractions}

\begin{annexitem}{annex:p10.dc.splitting}{thm:p10-dc-splitting}{Splitting and the exact inherited structure}
\begin{proof}
Put \(F=\operatorname{im}(f)\). Idempotence makes image and fixed-point
set equal. For directed \(A\subseteq F\), continuity gives
\(f(\bigsqcup A)=\bigsqcup f(A)=\bigsqcup A\), so the parent supremum
belongs to \(F\) and is its supremum there. Inclusion is therefore
continuous. The corestriction of \(f\) is continuous by the same
supremum equality, and \(rs=1_F\), \(sr=f\).

For a directed family \((f_i)\), the pointwise supremum \(f\) is
continuous by Theorem~\ref{thm:p10-domain-function-space}. Joint
continuity of composition gives
\(f^2(x)=\bigsqcup_{i,j}f_i(f_j(x))\).
For any \(i,j\) choose \(k\) above both in the directed family.
Monotonicity yields \(f_i(f_j(x))\sqsubseteq f_k(f_k(x))=f_k(x)\).
Conversely diagonal terms are among the terms, so this supremum is
\(\bigsqcup_k f_k(x)=f(x)\). Thus \(f^2=f\).

In any continuous retraction \((s,r)\), choose \(b\in B\) with
\(b\ll s(x)\). If \(x\sqsubseteq\bigsqcup A\) for directed
\(A\subseteq D\), continuity of \(s\) and way-below yield some
\(a\in A\) with \(b\sqsubseteq s(a)\), hence \(r(b)\sqsubseteq a\).
Thus \(r(b)\ll x\). The directed set \(r(B_{s(x)})\) has supremum
\(r(s(x))=x\), so \(r(B)\) is a basis. Countability passes to its
image. This argument does not assert \(r(b)\ll r(b')\) for arbitrary
\(b\ll b'\).

For the split image, each parent Scott-open restricts to an open because
inclusion is continuous. Conversely, if \(U\) is Scott-open in the
image, \(r^{-1}(U)\) is parent Scott-open and
\(s^{-1}(r^{-1}(U))=U\). Hence every image open is a restricted parent
open. This proves the subspace assertion without assuming that the
parent way-below relation restricts to that of every retract.
\end{proof}
\end{annexitem}

\begin{annexitem}{annex:p10.dc.algebraic-cover}{thm:p10-dc-algebraic-cover}{Two different ideal completions must not be conflated}
\begin{proof}
Use the reflexive restricted order on \(B\). Its ideals are nonempty
directed lower sets. The already proved ideal-completion theorem gives
an algebraic dcpo with directed unions and compact principal ideals.
For \(x\in D\), \(B_x\) is nonempty and directed by the basis theorem.
If \(a\sqsubseteq b\ll x\) with \(a,b\in B\), then \(a\ll x\);
thus it is also lower for the restricted order. Therefore \(s\) is
well defined and monotone. Basis interpolation gives
\(B_{\bigsqcup A}=\bigcup_{a\in A}B_a\) for directed \(A\), proving
continuity of \(s\).

Every ideal \(J\) is a directed subset of \(D\), so \(r(J)\) exists.
For a directed family \((J_i)\), its union is directed; each element
of that union lies in some \(J_i\). The defining upper-bound properties
give \(\bigsqcup\bigcup_iJ_i=\bigsqcup_i\bigsqcup J_i\), proving
continuity of \(r\). Finally \(r(s(x))=\bigsqcup B_x=x\).
Indeed this pair even has \(sr\sqsubseteq1\): if
\(b\ll\bigsqcup J\), some \(j\in J\) has \(b\sqsubseteq j\),
and lower closure puts \(b\) in \(J\).

The principal ideals of a countable \(B\) give a countable compact
basis of the cover. For the rational basis of \([0,1]\), the target's
only compact point is zero, while the cover has all principal rational
ideals compact. The target is not the directed supremum of its compact
points at any positive value. Hence this explicit retraction disproves
preservation of algebraicity, while preserving continuity.
\end{proof}
\end{annexitem}

\begin{annexitem}{annex:p10.dc.function-retract}{thm:p10-dc-function-retract}{Transport on the actual hom dcpo}
\begin{proof}
All displayed composites are continuous. For directed \((h_i)\),
pointwise suprema and continuity of \(s_E\) give
\(S(\bigsqcup_i h_i)(x)=\bigsqcup_i s_Eh_i r_D(x)\);
the identical argument with \(r_E\) proves continuity of \(R\).
The retraction equations give
\(R(S(h))=r_Es_Eh r_Ds_D=h\).
If both original pairs are embedding--projection pairs, put
\(q_E=s_Er_E\sqsubseteq1\), \(q_D=s_Dr_D\sqsubseteq1\).
For every \(k\), monotonicity gives
\(S(R(k))=q_Ekq_D\sqsubseteq k\), proving the second pair's
projection inequality. Coordinatewise products give continuous section
and retraction maps with the same equations and inequalities.
No enumeration or algorithm has been used, so this is not a proof that
the represented functions or representation maps are computable.
\end{proof}
\end{annexitem}

\begin{annexitem}{annex:p10.dc.plotkin-object}{thm:p10-dc-plotkin-object}{Finite matching and interpolation build the domain}
\begin{proof}
Transitivity of \(\ll_{\mathrm{EM}}\) follows by composing the
witnesses for each of its two quantified clauses.
Suppose finitely many nonempty finite sets \(M_j\) satisfy
\(M_j\ll_{\mathrm{EM}}N\). For each \(n\in N\), choose a predecessor
\(a_{j,n}\in M_j\) with \(a_{j,n}\ll n\). Also assign every
\(a\in M_j\) some \(n\in N\) with \(a\ll n\). For each \(n\),
these are a finite collection of approximants to \(n\). Interpolation
in the chosen basis supplies \(c_n\in B\) with each of those
predecessors way below \(c_n\), and \(c_n\ll n\).
Let \(C=\{c_n:n\in N\}\). Every member of each \(M_j\) is matched
to a \(c_n\); every \(c_n\) is matched from each \(M_j\), using
\(a_{j,n}\). Thus \(M_j\ll_{\mathrm{EM}}C\ll_{\mathrm{EM}}N\).
If the family of \(M_j\)'s is empty, choose any \(c_n\ll n\) in
the basis; this handles the empty interpolation premise too.
The abstract-basis and ideal-completion theorem now gives the continuous
object, its principal round-ideal basis and directed unions.

For \(x\in D\), the set \(\eta_D(x)\) is nonempty by its nonempty
basis approximants. It is lower for Egli--Milner approximation.
Given its finite sets \(M,N\), interpolate all their members together
to a single \(c\in B\) with every member way below \(c\ll x\).
Then \(M,N\ll_{\mathrm{EM}}\{c\}\), so it is rounded and directed.
It is monotone in \(x\). If \(M\in\eta_D(\bigsqcup A)\) for
directed \(A\), basis interpolation followed by the way-below definition
gives, for each \(b\in M\), an \(a_b\in A\) with \(b\ll a_b\).
One member of \(A\) lies above all the finitely many \(a_b\)'s.
Therefore \(M\) already belongs to some \(\eta_D(a)\). This proves
\(\eta_D(\bigsqcup A)=\bigcup_{a\in A}\eta_D(a)\).

For an algebraic basis, \(\ll\) on \(K(D)\) is just \(\sqsubseteq\).
Hence the finite-set relation is reflexive and its ideal completion is
algebraic. In a reflexive preorder, every ideal is the directed union
of its principal ideals; a compact ideal is below one of those contained
in it, hence equals it. Equal principal ideals are precisely the
mutually related finite sets. The finite subsets of a countable set are
countable. Finally \(\bot\ll x\) for every \(x\), and
\(L\ll_{\mathrm{EM}}\{\bot\}\) forces \(L=\{\bot\}\).
This principal ideal is contained in every ideal by lower closure and
is exactly \(\eta_D(\bot)\).
\end{proof}
\end{annexitem}

\begin{annexitem}{annex:p10.dc.formal-union}{thm:p10-dc-formal-union}{Ideal generation, not set-theoretic union of ideals}
\begin{proof}
Write \(M\prec N\) for \(\ll_{\mathrm{EM}}\) in this proof.
Finite set union preserves this relation in each pair of coordinates:
\(M\prec M',N\prec N'\) implies \(M\cup N\prec M'\cup N'\),
by retaining the witnesses in the corresponding sets.
The generators \(M\cup N\), for \(M\in I,N\in J\), form a
rounded directed family in this sense. Their lower closure is nonempty
by abstract-basis interpolation, lower by transitivity, and rounded
directed by enlarging the finitely many \(M\)'s inside \(I\) and
\(N\)'s inside \(J\). It contains the generators themselves because
each \(M,N\) has a successor in the relation inside its ideal.
This is exactly the displayed definition of formal union.

Commutativity follows from finite set union. Both bracketings of a
triple formal union are the lower ideal generated by
\(M\cup N\cup L\), for \(M\in I,N\in J,L\in K\): one inclusion
uses relation-preservation under union, and the reverse uses the fact
that each intermediate generator \(M\cup N\) belongs to \(I\cup J\).
Thus associativity holds. If \(M,N\in I\), rounded directedness
supplies \(C\in I\) with \(M,N\prec C\); then \(M\cup N\prec C\).
Lower closure proves \(I\cup I\subseteq I\). For \(L\in I\),
choose \(M\in I\) with \(L\prec M=M\cup M\), giving the reverse
inclusion. This proves idempotence.

For a directed family \((I_t,J_t)\) in the product of ideal domains,
membership in the formal union of the two coordinate suprema has
finite witnesses \(M\in I_t,N\in J_u\). A common upper index puts
both in one pair. This proves equality with the union of the pairwise
formal unions; monotonicity is immediate. Hence the operation is jointly
Scott-continuous, not merely separately continuous.

Finally the finite formal union of the \(\eta_D(b)\), \(b\in M\),
is generated by \(\bigcup_{b\in M}N_b\), with \(N_b\in\eta_D(b)\).
Every such union is \(\prec M\), proving containment in \(i(M)\).
Conversely, for \(L\prec M\), assign each \(a\in L\) to a matched
\(b\in M\), and ensure each \(b\) has at least one matched
predecessor. Interpolate the finite predecessors at \(b\) to
\(c_b\ll b\), with all those predecessors way below \(c_b\).
Then \(N_b=\{c_b\}\in\eta_D(b)\), and
\(L\prec\bigcup_bN_b\). Thus \(L\) belongs to the finite formal
union. Equality follows.
\end{proof}
\end{annexitem}

\begin{annexitem}{annex:p10.dc.free-powerdomain}{thm:p10-dc-free-powerdomain}{Existence and uniqueness in every dcpo-semilattice}
\begin{proof}
Define \(\tau_g(M)=\bigcup_{b\in M}^{A}g(b)\). This finite
operation is independent of ordering and repetition by the three
semilattice equations. If \(M\ll_{\mathrm{EM}}N\), choose finitely
many pairs \((a,b)\in M\times N\), with \(a\ll b\), covering every
element of both sets. Repeating entries along these pairs does not
change \(\tau_g(M)\) or \(\tau_g(N)\). Monotonicity of \(g\) and
of the finite operation therefore gives
\(\tau_g(M)\sqsubseteq\tau_g(N)\). Both Egli--Milner clauses are
needed here: no inflationary law \(x\sqsubseteq x\cup_A y\) is
available. Rounded directedness of an ideal consequently makes
\(\tau_g(I)\) directed, so the asserted supremum defines a map.

The map is monotone. For a directed family of ideals, an index
\(M\) belongs to their union exactly when it belongs to one member.
The two iterated directed suprema therefore coincide; this proves
Scott continuity of \(\operatorname{ext}(g)\).
For the unit, each \(M\in\eta_D(x)\) has \(g(b)\sqsubseteq g(x)\)
for every \(b\in M\), and idempotence gives
\(\tau_g(M)\sqsubseteq g(x)\). Conversely all singletons
\(\{b\}\), \(b\ll x\), occur in \(\eta_D(x)\). Their images
have supremum \(g(x)\) by continuity and the basis equation. Thus
\(\operatorname{ext}(g)\eta_D=g\).

For \(L\in I\cup J\), choose \(M\in I,N\in J\) with
\(L\ll_{\mathrm{EM}}M\cup N\). The matching argument bounds
\(\tau_g(L)\) by \(\tau_g(M)\cup_A\tau_g(N)\).
Conversely each generator \(M\cup N\) belongs to \(I\cup J\),
as proved above, so its value is included in the extension supremum.
Joint continuity of \(\cup_A\) gives
\[
 \operatorname{ext}(g)(I\cup J)
 =\bigsqcup_{M\in I,N\in J}
      (\tau_g(M)\cup_A\tau_g(N))
 =\operatorname{ext}(g)(I)\cup_A\operatorname{ext}(g)(J).
\]
Thus it is a homomorphism.

Every ideal \(I\) equals the directed union of \(i(M)\), \(M\in I\):
lower closure gives one inclusion and roundness gives the other.
The preceding theorem writes each \(i(M)\) as the finite formal
union of the \(\eta_D(b)\)'s. Any continuous homomorphism extending
\(g\) must consequently take \(i(M)\) to \(\tau_g(M)\), and \(I\)
to their directed supremum. This proves uniqueness, not merely leastness.

For a directed family \((g_j)\) in the function dcpo, finite joint
continuity gives \(\tau_{\bigsqcup_jg_j}(M)=\bigsqcup_j\tau_{g_j}(M)\).
All families involved are directed, and the pair-indexed family is
directed too. Interchanging its suprema yields
\(\operatorname{ext}(\bigsqcup_jg_j)(I)=
\bigsqcup_j\operatorname{ext}(g_j)(I)\). This proves the claimed
Scott continuity of extension. No approximation basis of \(A\) was used.
\end{proof}
\end{annexitem}

\begin{annexitem}{annex:p10.dc.powerdomain-functor}{thm:p10-dc-powerdomain-functor}{Universal comparisons prove the functor and monad laws}
\begin{proof}
The extension definition makes \(P^{\mathrm P}(u)\) a continuous
homomorphism with \(P^{\mathrm P}(u)\eta_D=\eta_Eu\).
The identity map on the powerdomain is a homomorphism with the same
generator values as \(P^{\mathrm P}(1)\), so uniqueness makes them
equal. For composable \(u,v\), the composite
\(P^{\mathrm P}(v)P^{\mathrm P}(u)\) and
\(P^{\mathrm P}(vu)\) are homomorphisms with generator values
\(\eta vu\); uniqueness proves composition. This equation on generators
also proves naturality of \(\eta\).
For a directed family \((u_j)\), continuity of \(\eta_E\) makes
composition \(u\mapsto\eta_Eu\) continuous; the preceding theorem's
continuity of extension therefore proves local continuity of \(P^{\mathrm P}\).
In particular it is monotone on hom orders. Retraction equations are
preserved by functoriality; projection inequalities are preserved by this
monotonicity and \(P^{\mathrm P}(1)=1\).

The extension of identity is well defined because the powerdomain is a
continuous domain and a dcpo-semilattice. Its defining equation gives
\(\mu_D\eta_{P^{\mathrm P}(D)}=1\).
Both \(\mu_DP^{\mathrm P}(\eta_D)\) and identity are homomorphisms
from the free algebra \(P^{\mathrm P}(D)\), and both take \(\eta_Dx\)
to \(\eta_Dx\), by naturality and the first unit equation. Uniqueness
gives the second unit equation.

For naturality of multiplication, compare
\(P^{\mathrm P}(u)\mu_D\) and
\(\mu_EP^{\mathrm P}(P^{\mathrm P}(u))\). Both are homomorphisms
from \(P^{\mathrm P}(P^{\mathrm P}(D))\), and their value on
\(\eta_{P^{\mathrm P}(D)}(I)\) is \(P^{\mathrm P}(u)(I)\).
Uniqueness makes them equal.
For associativity, the two displayed composites are homomorphisms from
\(P^{\mathrm P}(P^{\mathrm P}(P^{\mathrm P}(D)))\).
At a generator \(\eta_{P^{\mathrm P}(P^{\mathrm P}(D))}(I)\),
naturality and the first unit law make both values \(\mu_D(I)\).
Uniqueness proves equality. Thus all monad equations have been proved
through actual source and target maps, not inferred from the word
``flattening''.
\end{proof}
\end{annexitem}

\begin{annexitem}{annex:p10.dc.choice-example}{thm:p10-dc-choice-example}{Seven outcomes and the order that choice does not replace}
\begin{proof}
Every element of the finite flat domain is compact, so the Egli--Milner
relation is the two-sided lifting of its order. For nonempty finite
\(M,N\), write \(A_M=M\cap\{a,b\}\).
The first clause requires \(A_M\subseteq A_N\), since a nonbottom
element can only be matched above to itself. If \(\bot\notin M\),
the second clause forces \(\bot\notin N\) and
\(A_N\subseteq A_M\), so then only equality is possible.
If \(\bot\in M\), it matches every target; the second clause imposes
no extra condition. Mutual relatedness therefore forces equal sets:
when just one contains bottom, the reverse second clause fails; when
both contain bottom, the two first clauses force \(A_M=A_N\).
Thus the seven sets give seven distinct principal ideals.
In a finite preorder each ideal has an upper generator in itself,
obtained by finitely many directedness steps, and equals its principal
ideal. Hence these are all elements, not just seven selected examples.

The displayed comparisons follow immediately from the two clauses.
Formal union of principal ideals corresponds to actual union of their
finite generators, because the relation is reflexive and the union
operation preserves its order. In particular the union of the two
singleton outcomes is represented by \(\{a,b\}\).
But the principal ideal below \(\{a\}\) has only \(\{\bot\}\),
\(\{\bot,a\}\), \(\{a\}\), and the analogous ideal below \(\{b\}\)
has the corresponding three sets. Their ordinary union contains
\(\{a\},\{b\}\) but no upper bound for both within that union.
It is not an ideal. The formal union does contain the missing common
generator \(\{a,b\}\).

An empty source set cannot match each element of a nonempty target in
the second clause; a nonempty source cannot match any element into an
empty target in the first clause. Thus adding the empty finite set
creates an isolated class, unrelated to every nonempty class, and not
the existing bottom class. It would require a changed signature/theory
to represent a nullary empty choice in the desired way.
The constant map to \(1/2\) is continuous and idempotent, and its
one-point image is compact. The strictly increasing rational points
below \(1/2\) have supremum \(1/2\) without any reaching it, so that
point is not compact in the parent. This verifies the advertised
ordinary-retract boundary independently of the algebraic-cover example.
\end{proof}
\end{annexitem}
\begin{annexitem}{annex:p10.dc.effective-choice}{thm:p10-dc-effective-choice}{Recursive compact codes and enumerable approximation matrices}
\begin{proof}
Identity has the recursive matrix \((n,m)\mapsto(b_m\sqsubseteq b_n)\).
For effective \(f:D\to E\), \(g:E\to F\), with compact names
\(b_n,c_m,d_k\), algebraicity and continuity give
\[
 (n,k)\in R_{gf}\quad\Longleftrightarrow\quad
 (\exists m)\ (n,m)\in R_f\ \text{and}\ (m,k)\in R_g.
\]
For the forward implication, write \(f(b_n)\) as the directed supremum
of its compact approximants. Continuity of \(g\) makes \(g(f(b_n))\)
the directed supremum of their images. Compactness of \(d_k\) gives
a compact \(c_m\sqsubseteq f(b_n)\) with \(d_k\sqsubseteq g(c_m)\).
The reverse implication follows from monotonicity of \(g\).
The existential composition of two recursively enumerable relations is
recursively enumerable: enumerate both matrices in dovetailing stages and
output each outer pair when a common middle witness has appeared. Hence
effective maps compose; identities and associativity are their actual
continuous-map identities and composition.

For a computable \(x\), the identical compactness/continuity argument gives
\[
 I_{f(x)}=\{m:(\exists n)\ n\in I_x\ \text{and}\ (n,m)\in R_f\}.
\]
Dovetailing the two enumerations gives a procedure enumerating this set.
It may not decide membership; recursive enumerability, not decidability,
is the asserted property of a computable point.

Code the nonempty finite compact basis sets by finite index lists.
Every list denotes a principal ideal; the preceding algebraic theorem
shows these enumerate all compact points of the powerdomain. To compare
two such points, test both finite forall-exists Egli--Milner clauses with
the recursive compact order. This terminates. Mutual comparison handles
different lists naming the same compact point; choosing a canonical
representative or deciding equality of arbitrary domain points is unnecessary.

For finite compact sets \(M\subseteq K(D)\), \(N\subseteq K(E)\),
the matrix entry
\(i_E(N)\sqsubseteq P^{\mathrm P}(f)(i_D(M))\) holds precisely when
there are nonempty finite compact sets \(N_a\), one for each \(a\in M\),
such that every \(c\in N_a\) satisfies \(c\sqsubseteq f(a)\) and
\[
 N\sqsubseteq_{\mathrm{EM}}\bigcup_{a\in M}N_a.
\]
Indeed functoriality on generators and preservation of formal union give
\(P^{\mathrm P}(f)(i_D(M))=\bigcup_{a\in M}^{\mathrm{formal}}\eta_E(f(a))\).
The ideal-operation formula in the reflexive compact basis is exactly
the stated finite-witness condition. For each proposed list of witnesses,
its finitely many required entries in \(R_f\) can be recognized by
dovetailing its enumeration. The finite Egli--Milner comparison is
decidable. Dovetail over all lists of witnesses and output an input/output
pair when its tests succeed. This enumerates the matrix of \(P^{\mathrm P}(f)\),
proving effectivity without deciding every entry of \(R_f\).

The unit's compact matrix is even recursive:
\(i(M)\sqsubseteq\eta_D(b_n)\) iff every \(b_j\in M\) satisfies
\(b_j\sqsubseteq b_n\). A compact element of the twice-applied
powerdomain is generated by a finite list
\(i(M_1),\ldots,i(M_t)\) of compact elements of the first.
Its flattening is \(i(M_1\cup\cdots\cup M_t)\), since flattening
extends identity and preserves formal union. Thus its compact matrix
tests \(N\sqsubseteq_{\mathrm{EM}}M_1\cup\cdots\cup M_t\), also
recursively. The already established functor and monad equations are
equalities of these effective maps, so they hold in the restricted
category too. For the stated failure case, the powerset compact basis is
the effectively numbered finite subsets of the natural numbers, with
decidable inclusion. Every constant map is Scott-continuous. If its
constant value \(U\) made its compact approximation matrix recursively
enumerable, testing the compact singleton codes would enumerate exactly
the members of \(U\), contrary to the hypothesis. Hence this map is
not effective. Nothing in this proof gives a reduction semantics or an
operational adequacy theorem.
\end{proof}
\end{annexitem}
